% Generated draft: requires mathematical and source review.
\documentclass[11pt]{article}
\usepackage[margin=25mm]{geometry}
\usepackage{fontspec}
\setmainfont{Latin Modern Roman}
\usepackage{amsmath,amssymb,mathtools,unicode-math}
\setmathfont{Latin Modern Math}
\usepackage{enumitem,needspace,xurl,hyperref}
\setlength{\parindent}{0pt}
\setlength{\parskip}{5pt}
\setlength{\emergencystretch}{3em}
\newcommand{\N}{\mathbb{N}}
\newcommand{\Z}{\mathbb{Z}}
\newcommand{\Q}{\mathbb{Q}}
\newcommand{\R}{\mathbb{R}}
\newcommand{\C}{\mathbb{C}}
\newcommand{\F}{\mathbb{F}}
\newcommand{\A}{\mathbb{A}}
\newcommand{\PP}{\mathbb{P}}
\newcommand{\E}{\mathbb{E}}
\newcommand{\eps}{\varepsilon}
\newcommand{\dd}{\,\mathrm{d}}
\begin{document}
\section*{UnsolvedMath ID: 30001727}
\textbf{Source problem code:} catalogue entry
\subsection{Definitions and mathematical statement}
The supplied record poses a historiographical and conceptual research question, rather than a quantified mathematical conjecture. It supplies no definition of a \emph{mathematical table}, no specified class of \emph{digital tables}, and no criterion for deciding when a generalization undermines the notion of a table.

\begin{quote}
What exactly is meant by a mathematical table in the digital era, and would generalizing the notion enough to include digital tables undermine the notion itself?
\end{quote}

The research task is to clarify the scope of these terms and assess the conceptual consequences of extending the notion of a mathematical table to digital tables. The intended objects, admissible generalizations, and assessment criteria remain unspecified. Consequently, a precise mathematical definition or proposition cannot be extracted without adding assumptions absent from the source; the question is retained with these qualifications.
\subsection{Short English statement}
Clarify what counts as a mathematical table in the digital era and whether extending the concept to digital tables undermines its meaning. The source leaves the relevant definitions and assessment criteria unspecified, so this remains a conceptual research question requiring human review.
\subsection{Sources}
\begin{itemize}
\item Local UnsolvedMath registry, supplied record with permanent ID \(30001727\), title \emph{Scope of Digital Mathematical Tables}, source problem code \texttt{catalogue entry}, in the Oberwolfach Reports Open Problems collection. The recorded proposed year is 2011; no individual proposer is identified.
\item \emph{Mini-Workshop: History of Numerical and Graphical Tables} (2011), Oberwolfach Reports, pp.\ 639--689, DOI \texttt{10.4171/owr/2011/12}. Supplied source citation and literature reference, labeled \emph{Oberwolfach Report 12/2011}: \url{https://doi.org/10.4171/owr/2011/12}. The supplied background describes these reports as documenting mathematical workshops and their open-problem discussions.
\item Supplied statement-verification annotation: \texttt{corrected\_verified}. The annotation states that the official report confirms this historiographical research question and that an extraction label was removed. This is supplied verification metadata, not an independent verification performed for this draft.
\item Literature assessment dated exactly \texttt{2026-08-21}: the supplied curation labels the question \emph{open}, describes the scope and definition of mathematical tables in the digital era as unresolved, and reports that no definitive literature-wide resolution was found. This is supplied dated curation derived from the OWR material, not an independently verified literature assessment. The source explicitly calls for independent verification of time-sensitive status and cited claims.
\end{itemize}
\end{document}
